\chapter{File formats and data}
\label{ch:formats}

\section{Products}

\begin{center}
\begin{tabularx}{\textwidth}{@{}lX@{}}
\toprule
File & Content \\
\midrule
\file{name.fbk.md} & the report: summary, analysis sections, findings,
  provenance, references (complete citations + paste-ready BibTeX) \\
\file{name.fbk.json} & the same content as data: \code{title}, \code{sections},
  \code{findings} (severity, message, rule id, suggested fix, refusals),
  \code{provenance}, plus a few program-specific keys \\
\file{structure.fbk.inp} & a generated ORCA input (pure ASCII) \\
\file{input.fix\_\textit{variant}.inp} & a corrected input from \menu{9} \\
\bottomrule
\end{tabularx}
\end{center}

Nothing carries timestamps; products are deterministic.

\section{Input JSON: cross-level records (menu 7)}

\begin{codeblock}
[
  {"label": "No.1",
   "occupation": "(5f)^1 ...",        # optional, informational
   "energies": {"HF": -1930.87692, "CCSD(T)": -1932.97951}},
  {"label": "No.2", "energies": {"HF": -1930.87392, "CCSD(T)": -1932.97745}}
]
\end{codeblock}

The top level may also be an object with a \code{records} list. All records
must carry the same level names; labels must be unique.

\section{Input JSON: crystal-field fit (menu 10)}

\begin{codeblock}
{
  "point_group": "C3",                  # C3 | Oh | C1/none
  "J": 8.0,                             # the |J M> manifold
  "levels": [0.0, 12.3, ...],           # one energy per sampled state
  "coefficients": [[[re, im], ...], ...],  # amplitudes on |J M>, M = m - J
  "declaration": {                      # optional: the five A5 items
    "convention": "...", "projection": "...", "units": "cm^-1",
    "z_axis": "...", "origin": "..."},
  "compare": {                          # optional: a second parameter set
    "label": "...",
    "parameters": {"2,0": 1879.0, "2,2": -43.0},
    "declaration": {"projection": "..."}}
}
\end{codeblock}

Parameter keys are \code{"k,q"} strings or two-element lists. Pass the
amplitudes themselves, not their complex conjugates (the conjugation trap,
Chapter~\ref{ch:menu10}).

\section{Input JSON: magnetic-doublet table (menu 16)}

\begin{codeblock}
{
  "system": "...",                       # optional label
  "reference": 0,                        # optional: index or label
  "doublets": [
    {"label": "ground", "g": [0.41, 0.44, 9.04], "energy": 0.0,
     "axis3": [0.0, 0.0, 1.0]},          # unit axis of the largest g value
    {"label": "1st excited", "g": [0.41, 0.44, 9.04], "theta3": 9.52}
  ]
}
\end{codeblock}

Per doublet, \code{theta3} (degrees) or \code{axis3} says where the largest
$g$ axis points; the reference doublet may omit both (its \code{theta3} is 0
by definition when the other angles are computed from axes, and it must carry
\code{axis3} then). \code{energy} is optional and printed when present.

\section{Input JSON: cross-structure mapping manifest (menu 17)}

\begin{codeblock}
{
  "structures": [                         # one entry per structure, same basis
    {"name": "r=1.094", "export": "a.json"},
    {"name": "r=2.600", "export": "b.json"}
  ],
  "tau": 0.5,                             # optional: one threshold for both criteria
  "selections": {"r=1.094": [4, 5]},      # optional: orbitals to make consistent
  "active": {"r=1.094": [4, 5, 6],        # optional: active-space overlap check
             "r=2.600": [4, 5, 6]}        # (per adjacent pair, in list order)
}
\end{codeblock}

Export paths are relative to the manifest; \code{selections} and \code{active}
name 0-based orbital indices per structure name.  Without \code{tau} the
threshold is read from the non-matchable curve (the source's Eq.~(8)).

\section{Input JSON: guess-transfer manifest (menu 18)}

\begin{codeblock}
{
  "structures": [                         # the neighbours, same basis/atom order
    {"name": "r=1.094", "export": "a.json"},
    {"name": "r=2.600", "export": "b.json"}
  ],
  "template": {                           # the target geometry
    "export": "target.json",              # its metric (S-Matrix)
    "mkl": "target.mkl"},                 # orca_2mkl <target> -mkl
  "delta": 1.0                            # optional neighbourhood radius (Angstrom)
}
\end{codeblock}

The output \code{<template stem>.fbk.mkl} converts with
\code{orca\_2mkl <template stem>.fbk -gbw} and is read with \code{!moread} +
\code{\%moinp}.  The mkl is a plain-text Molekel file (geometry in \AA,
verbatim \code{\$BASIS}, grouped coefficients, occupations); the toolkit's
reader/writer translates its measured p-shell row order to and from the
export convention.

\section{Input JSON: DM-AS selection manifest (menu 20)}

\begin{codeblock}
{
  "candidates": [                         # one entry per scanned space
    {"nel": 6, "norb": 8, "output": "cand_e6o8.out"}
  ],
  "reference": {"output": "dft.out"},     # the SCF block of that output is used
  "protocol": "gdm"                       # optional: gdm (default) | vgdm
}
\end{codeblock}

The reference may instead be a supplied value,
\code{\{"magnitude\_debye": 1.85, "source": "NIST (gas phase)"\}}, optionally
with a \code{vector} for the directional protocol.  Candidate outputs are
read at their single-root dipole block (State: 0); a state-averaged output is
refused.  Menu~\menu{19} writes this manifest (with the candidate list
pre-filled) next to the batch.

\section{The orca\_2json export and its request file (menus 12--15, 17, 21)}

\code{orca\_2json <base>.gbw} rewrites an orbital file into
\code{<base>.json}; which blocks it carries depends on the request file
\code{<base>.json.conf} placed next to the \code{.gbw}.  The requests the
menus use:

\begin{center}
\begin{tabularx}{\textwidth}{@{}>{\raggedright\arraybackslash}p{4.6cm} X@{}}
\toprule
Request & Blocks and measured conventions \\
\midrule
\code{\{"MOCoefficients": true\}} & Orbitals, occupations, energies; the base of
  every export. \\
\code{"1elIntegrals": ["H", "S", "T", "V"]} & \code{S-}/\code{H-}/\code{T}-/\code{V-Matrix};
  the cross-structure mapping (menu~\menu{17}) reads the kinetic matrix. \\
\code{"FockMatrix": ["J", "K"]} & Density-derived Fock blocks, one square
  matrix per spin; menu~\menu{21}.  Measured: the exported \code{K} block is
  the exchange contribution as it enters the Fock, so the source's
  $\tfrac12 K_{aa}$ equals $-\operatorname{diag}(CKC^{\mathsf T})_a =
  \sum_i (ai|ai)$; the exported \code{F} block equals $J+K$ and \emph{excludes}
  the core Hamiltonian ($H$ must be requested too for the full Fock). \\
\code{"2elIntegrals": ["MO\_IAJB"]},\newline
\code{"OrbWin": [i0, i1, a0, a1, 0, 0, 0, 0]} & Windowed two-electron
  entries \code{[i, j, a, b, value]} with \code{value} the chemist-notation
  $(ia|jb)$ and the internal indices first; menu~\menu{21} (\code{apcx}).  The
  input window is \emph{eight} integers (the second window all zeros for one
  window; the four-integer form is rejected by \code{orca\_2json}), inclusive
  0-based; the record written back into the JSON is the four-integer form. \\
\bottomrule
\end{tabularx}
\end{center}

Measured on ORCA 6.1.1: a \code{FockMatrix} request is answered from the run's
\code{<base>.densities} and \code{<base>.densitiesinfo} files, so the export
must be produced where the run wrote (or with the whole file family moved
together, names kept) -- a bare \code{.gbw} copy fails with ``NO densities are
available''.

\section{Structure and charge input (menus 2 and 11)}

Standard XYZ (atom count, comment, \code{element x y z}). For \menu{11}, the
interactive prompts take per-element charges (\code{O=-2,H=0.4}) and optional
radial moments \code{r2,r4,r6} in \AA$^k$. An element present in the structure
but missing from the charge list is refused, with the missing symbols named.

\section{Knowledge data (inside the package)}

\begin{center}
\begin{tabularx}{\textwidth}{@{}lX@{}}
\toprule
File & Content \\
\midrule
\file{knowledge/sources.bib} & the reference database (BibTeX). Every
  literature-evidence item must carry a bibkey that resolves here; loading
  fails otherwise \\
\file{knowledge/rules/*.yaml} & the rule base: id, kind (recipe/diagnosis),
  condition tree, action, refusals, evidence, confidence, severity (diagnosis
  only) \\
\file{knowledge/basis\_ecp.yaml} & the basis/ECP recommendation table \\
\file{knowledge/templates/*.j2} & the input templates (ORCA, OpenMolcas) \\
\file{knowledge/toolindex/tools.yaml} & the external-tool index (relation and
  status per entry) \\
\file{ui/menu.yaml} & the menu definition: number, title, handler (append-only
  numbers) \\
\bottomrule
\end{tabularx}
\end{center}

\section{The report as an interface}

The Markdown report is for humans; the JSON companion is the scripting
interface. Its stability follows the release policy: within a minor version
the section titles and finding fields are stable; new fields may be added.
Finding \code{rule\_id}s are stable identifiers you may key on.
